% Vai trò nghiên cứu trong Long-CoT (A1--A4).
% Bản tiếng Việt: 2026-11-08; giữ terminology và nguồn của bản trước.

\section{Các vai trò nghiên cứu trong Long-CoT}
\label{sec:research_roles}

Để tổ chức literature về suy luận Long Chain-of-Thought (Long-CoT), chúng tôi chia các hướng nghiên cứu thành bốn trục kỹ thuật chính:
\begin{enumerate}
    \item \textbf{A1: Training/Data}---mục tiêu tối ưu, quy trình reinforcement learning (RL), tuyển chọn dữ liệu theo curriculum và các pipeline distillation;
    \item \textbf{A2: Inference/Control}---early stopping thích ứng, điều hướng cây tìm kiếm, thoát bẫy suy luận và cô lập context theo mô-đun;
    \item \textbf{A3: Analysis}---mô hình hóa cấu trúc diễn ngôn, kiểm tra nhân quả bằng can thiệp phản thực, giới hạn scaling theo độ phức tạp tổ hợp và hành vi overthinking;
    \item \textbf{A4: Evaluation/Benchmark}---bài toán mang trạng thái qua horizon dài, chỉ số redundancy dựa trên graph và giao thức process supervision.
\end{enumerate}
Các vai trò này là góc nhìn về đóng góp kỹ thuật, không phải những loại model loại trừ nhau. Một công trình có thể liên quan nhiều vai trò; chúng tôi gán \emph{primary role} theo đóng góp trung tâm và dùng \emph{secondary roles} để ghi nhận phần giao thoa. Các subgroup là cách tổ chức được nhóm đề xuất từ những phương pháp đã có, nhằm làm rõ các hướng nghiên cứu khác nhau.

\paragraph{Phạm vi và quy tắc gán vai trò.}
Long-CoT ở đây chỉ quá trình deliberation mở rộng, có exploration hoặc verification; không áp dụng một ngưỡng token chung. Chỉ có đầu vào long-context chưa đủ để thuộc phạm vi này. A1 thay đổi dữ liệu training hoặc tham số model; A2 thay đổi cách thực thi inference; A3 chủ yếu khảo sát một hiện tượng; A4 chủ yếu cung cấp công cụ hoặc giao thức đánh giá. Danh mục gồm 24 công trình, phân bố 6/5/8/5 theo primary role: 19 nghiên cứu trực tiếp Long-CoT, bốn công trình suy luận liên quan và một công trình nền tảng về phương pháp kiểm tra faithfulness. Dẫn chéo không được tính thành paper bổ sung. Đây là bản tổng hợp chọn lọc, không phải systematic review đầy đủ; kết quả được báo cáo chưa được nhóm tái lập độc lập.

% -------------------------------------------------------------------------
\subsection{Training/Data (A1)}
\label{subsec:a1_training_data}

A1 nghiên cứu cách hình thành năng lực deliberation qua horizon dài trong post-training, từ pure RL đến giám sát dựa trên tập dữ liệu nhỏ được chọn kỹ và pipeline curriculum nhiều giai đoạn.

\subsubsection{Hình thành và mở rộng Long-CoT bằng RL}
Một hướng quan trọng khảo sát khả năng hình thành hành vi suy luận phức tạp trực tiếp từ RL quy mô lớn với rule-based verification. Trong \textbf{DeepSeek-R1-Zero} \cite{rra1_deepseek_r1}, policy được tối ưu trực tiếp từ \texttt{DeepSeek-V3-Base} bằng Group Relative Policy Optimization (GRPO). Thuật toán tính advantage trong nhóm output được lấy mẫu mà không cần critic network riêng:
\begin{equation}
    A_i = \frac{r_i - \text{mean}(\{r_j\})}{\text{std}(\{r_j\}) + \epsilon},
\end{equation}
trong đó reward $r_i$ phản ánh độ đúng đáp án và định dạng tag. Trong thiết lập này, DeepSeek-R1-Zero tự hình thành self-verification và các trajectory exploration mở rộng, đạt 77.9\% Pass@1 trên AIME 2024 và 95.9\% trên MATH-500. Tuy nhiên, tác giả ghi nhận hiện tượng trộn ngôn ngữ và khó đọc. Để xử lý, \textbf{DeepSeek-R1} hoàn chỉnh sử dụng cold-start SFT với vài nghìn trace được chọn lọc, reasoning RL, rejection sampling tạo tập dữ liệu tổng hợp 800k mẫu và RL căn chỉnh đa tác vụ. Model đạt 79.8\% trên AIME 2024 (671B MoE); SFT distillation cũng chuyển giao hiệu quả sang dense model nhỏ hơn, với 72.6\% trên Qwen-32B.

\textbf{Kimi k1.5} \cite{rra1_kimi_k15} mở rộng context training đến 128k token cho cả văn bản và thị giác. Phương pháp dùng Online Policy Mirror Descent với relative entropy regularization, policy không có value head và gradient âm từ rollout sai. Để chấm các trajectory dài, tác giả huấn luyện Chain-of-Thought Reward Model (CoT-RM) trên 800k mẫu: model sinh rationale trước khi phán quyết độ đúng đáp án cuối ở dạng JSON. Accuracy trên mẫu spot-check thủ công do tác giả báo cáo khoảng 98.5\%, so với khoảng 84.4\% của outcome reward model truyền thống; Kimi k1.5 đạt 77.5\% Pass@1 trên AIME 2024. Khung \emph{Long2short} dùng shortest-path rejection sampling và preference optimization có length penalty để chuyển tri thức từ suy luận dài sang trace ngắn hơn, nhằm giảm latency lúc inference.

\subsubsection{SFT với tập dữ liệu nhỏ được chọn lọc}
Khác với RL hoặc distillation corpus quy mô lớn, hướng data-centric khảo sát liệu một tập demonstration nhỏ nhưng chất lượng cao có đủ để kích hoạt năng lực suy luận trên base model mạnh. \textbf{s1} \cite{rra1_s1} xây dựng \texttt{s1K}, gồm 1,000 bài khó với reasoning trace được distill từ \texttt{gemini-2.0-flash-thinking-exp-1219}.

Fine-tune \texttt{Qwen2.5-32B-Instruct} trên s1K đạt 50.0\% Pass@1 trên AIME 2024 với greedy decoding thông thường. Khi kết hợp \emph{budget forcing}---chặn delimiter \texttt{</thought>} và nối thêm ``Wait''---kết quả đạt 56.7\% trên AIME 2024 và 59.6\% trên GPQA Diamond. Đây là bằng chứng về sample efficiency của tập dữ liệu được chọn lọc, trong điều kiện năng lực sẵn có của base model.

\textbf{LIMO} \cite{rra1_limo} chọn 817 mẫu theo độ khó câu hỏi và rubric bốn chiều của trace: lập luận chi tiết, self-verification, exploration và độ hạt thích ứng. SFT trên \texttt{Qwen2.5-32B-Instruct} đạt 63.3\% Pass@1 trên AIME 2024 và 95.6\% trên MATH-500. Tuy nhiên, khi chuyển cùng quy trình tối giản sang \texttt{Qwen1.5-32B-Chat}, accuracy trên AIME giảm còn 9.2\%, thấp hơn 54.1 điểm phần trăm. Kết quả cho thấy hiệu quả dữ liệu phụ thuộc chặt vào base model cụ thể được thử nghiệm.

\subsubsection{Pipeline curriculum nhiều giai đoạn}
\textbf{Light-R1} \cite{rra1_light_r1} công bố quy trình có thể tái lập gồm ba stage training và model merging, khởi tạo từ \texttt{Qwen2.5-32B-Instruct} thay vì pretraining từ đầu. Pipeline gồm Stage~1 SFT trên 76k bài toán đã lọc với sequence length tối đa 20k token; Stage~2 SFT trên 3,000 bài thi đấu; semi-on-policy DPO với context 32k; và merging trong không gian trọng số. Accuracy của Light-R1-32B trên AIME 2024 tăng từ 16.6\% (base) lên 69.0\% (+SFT-1), 73.0\% (+SFT-2), 75.8\% (+DPO) và 76.6\% (+Merging). GPQA lại biến động: 64.3\% ở SFT-1, 60.6\% ở SFT-2 và 61.8\% sau merging. Tác giả còn kiểm tra decontamination và ghi nhận test-set leakage trong một số corpus reasoning mở trước đó.

\subsubsection{RL hướng efficiency và học phân bổ compute}
Arora và Zanette~\cite{rra1_efficient_rl} tiếp tục training reasoning model bằng reward độ đúng kết hợp length penalty có thể điều chỉnh. Với distilled model 7B tại $\alpha=0.2$, mức tiết kiệm token khác nhau rõ theo tác vụ: 83\% trên GSM8K và 27\% trên AIME 2024. Đây là đánh đổi efficiency--accuracy, không phải bảo đảm giữ nguyên accuracy trong các lần training hữu hạn. Compression cũng có thể làm giảm faithfulness đo được. Hướng này bổ sung cho RL tập trung vào capability bằng cách học khi nào deliberation mở rộng xứng đáng với chi phí.

% -------------------------------------------------------------------------
\subsection{Inference/Control (A2)}
\label{subsec:a2_inference_control}

Các phương pháp inference control điều tiết compute budget trong generation, xử lý overthinking, sự bám giữ lỗi sớm và overhead của context qua horizon dài.

\subsubsection{Early stopping thích ứng và đánh giá sufficiency}
Long-CoT thường tiếp tục verification và thử cách giải khác sau khi đã đủ cơ sở kết luận. \textbf{DTSR} \cite{rra2_dtsr} xử lý mà không retraining: theo dõi reflection trigger như ``Wait'', ``Alternative'' để chọn checkpoint, rồi dùng metacognitive sufficiency check đánh giá phần suy luận hiện tại đã đủ để chốt đáp án chưa. Trên \texttt{Qwen3-14B}, DTSR giảm token từ 5,761 xuống 3,748 ($-34.9\%$), với accuracy 84.8\% so với 84.4\% của baseline. Tuy nhiên, trên \texttt{Qwen3-8B}, accuracy giảm 0.9 điểm phần trăm, từ 81.9\% xuống 81.0\%, đổi lấy mức giảm 32.0\% token. Khả năng tự đánh giá sufficiency vì thế thay đổi theo quy mô model.

\textbf{REFRAIN} \cite{rra2_refrain} xây dựng điều kiện dừng từ reflective redundancy. Bộ phân biệt kiểm tra giả thuyết đáp án đã có, reflection token và sentence embedding similarity vượt ngưỡng $\tau$; ngưỡng được điều chỉnh bằng Sliding-Window UCB bandit. Trên \texttt{gpt-oss-20B}, REFRAIN cải thiện Pass@1 ở các benchmark được thử: +2.73 điểm phần trăm GSM8K, +3.40 điểm phần trăm MATH-500 và +7.07 điểm phần trăm GPQA-Diamond, đồng thời giảm 33.3\% đến 44.7\% token.

\subsubsection{Phát hiện bẫy và phục hồi thích ứng}
Trong suy luận dài, một cam kết sai sớm có thể bám dai trong prefix, khiến self-correction ở cuối vẫn bị neo vào tiền đề sai. \textbf{TAAR} \cite{rra2_taar} ghi nhận 44.4\% thinking trap được chẩn đoán xuất hiện trong 20\% đầu sequence. Phương pháp huấn luyện diagnostic policy $\pi_\phi$ dự đoán vị trí bẫy và xác suất thoát; \emph{Cut@Trap} cắt trajectory trước đoạn bẫy rồi restart với temperature thích ứng. Với cùng path budget được báo cáo $K=4$ trên \texttt{DeepSeek-R1-Distill-Qwen-8B}, accuracy trung bình trên năm benchmark thi đấu tăng từ 64.5\% (Avg@4) và 66.0\% (PRM@4) lên 69.1\%; Cut@Trap có escape rate 16.5\%. Trên \texttt{GPT-OSS-120B}, kết quả lại giảm nhẹ từ 80.6\% xuống 79.2\%: model mạnh gặp ít bẫy hơn, trong khi pruning mạnh có nguy cơ bỏ prefix gần đúng.

Ở độ hạt lớn hơn, \textbf{MCTSr} \cite{rra2_mctsr} kết hợp Monte Carlo Tree Search với self-refinement bằng ngôn ngữ tự nhiên; mỗi node là một lời giải toán hoàn chỉnh. Dùng quy tắc UCT được sửa đổi, phương pháp sinh critique, viết lại lời giải và cập nhật value qua self-evaluation. Trên \texttt{LLaMA-3-8B}, search với 8 rollout tăng accuracy MATH từ 24.36\% (zero-shot) lên 58.24\%; ở Level~5, kết quả tăng từ 7.10\% lên 34.06\%.

\subsubsection{Cô lập context và quản lý trạng thái}
Autoregressive generation trên một flat context có attention cost bậc hai $O(T^2)$ và nguy cơ pha loãng context khi sequence dài. \textbf{Recursive Models} \cite{rra2_recursive_models} tổ chức subtask trong các stack frame độc lập bằng tag \texttt{<call>} và \texttt{<return>}. Khi subtask hoàn thành, chỉ return value ngắn được chuyển về caller; token trung gian được loại khỏi active context. Trên SAT, model đệ quy fine-tune từ \texttt{Qwen2.5-3B-Instruct} đạt 98\% (Easy), 95\% (Medium) và 64\% (Hard OOD), vượt flat CoT baseline. Trên sinh trace cờ Go 4x4 với length-OOD, model giữ 38.5\% trace accuracy, trong khi flat CoT còn 1.0\%.

\subsubsection{Budget forcing và điều khiển độ sâu exploration}
Hai cơ chế giao thoa bổ sung cho early stopping: budget forcing của s1~\cite{rra1_s1} điều tiết độ dài deliberation, còn TIP~\cite{rra3_underthinking} hạn chế đổi hướng suy nghĩ quá sớm bằng penalty lên logits của transition token trong một cửa sổ decoding. Chúng nhắm vào hai lỗi phân bổ khác nhau: tiếp tục sau khi đã đủ lời giải và bỏ một nhánh triển vọng trước khi khám phá đủ sâu. Các bản ghi primary role vẫn lần lượt nằm ở A1 và A3.

% -------------------------------------------------------------------------
\subsection{Analysis (A3)}
\label{subsec:a3_analysis}

Các nghiên cứu phân tích khảo sát hành vi suy luận, từ cấu trúc diễn ngôn và quan hệ nhân quả nội tại đến graph mô tả failure mode và giới hạn độ phức tạp tổ hợp.

\subsubsection{Cấu trúc diễn ngôn và quan hệ nhân quả nội tại}
Một câu hỏi nền tảng là quan hệ diễn ngôn nhìn thấy trong trace có phản ánh dependency của tính toán nội tại hay không. \textbf{ReasoningFlow} \cite{rra3_reasoningflow} biểu diễn trace bằng Directed Acyclic Graph (DAG), với các cạnh định kiểu \emph{support, attack, verify, assumption}. Để kiểm tra dependency nhân quả, tác giả dùng \emph{Thought Anchors}, đo log KL-divergence của dự đoán bước $s_j$ khi attention đến bước $s_i$ bị mask. Trên QwQ và AIME, điểm Thought Anchors dự đoán cạnh ReasoningFlow có precision-recall tương đương heuristic chọn node gần theo khoảng cách token; mức alignment vì thế hạn chế trong thiết lập đã thử. Trong phân tích hành vi trace (Hình~5), local verification kèm sửa lỗi có step accuracy 84\%, so với 64\% ở bước lỗi trước đó; bước đã sửa có tỷ lệ dẫn đến đáp án cuối đúng 20\% so với 5\%. Đây là liên hệ quan sát trên trace, không phải can thiệp nhân quả độc lập có đối chứng.

\textbf{Lanham và cộng sự} \cite{rra3_lanham_faithfulness} khảo sát faithfulness trên model đến 175B tham số bằng early answering, chèn lỗi và filler token. Trong \S2.5.1 và Hình~5, thay CoT bằng filler không mang thông tin (``\texttt{ ...}'') không tăng task accuracy trên các benchmark được đánh giá. Accuracy giảm nhẹ ở TruthfulQA và OpenBookQA; tác giả nêu out-of-distribution effect như một cách giải thích có thể. Trong các model và tác vụ này, chỉ tăng test-time token mà không có nội dung ngữ nghĩa hữu ích không đem lại lợi ích reasoning; không thể khái quát thành quy luật cho mọi kiến trúc latent-compute.

\subsubsection{Biểu diễn graph và chẩn đoán failure mode}
Vượt ra ngoài chỉ số sequence tuyến tính, \textbf{CoT2Graph} \cite{rra3_cot2graph} dùng typed dependency graph và Graph Neural Network (GNN), kết hợp GNNExplainer để nhận diện dấu hiệu cấu trúc. Tác giả quan sát trace đúng thường dựa vào strategic planning sớm, còn trace sai có synthesis node quá tải và dependency chưa được đáp ứng. Attribution này giải thích GNN predictor, không phải cơ chế nhân quả của LLM. Trên GPQA với trace \texttt{DS-Qw14B}, accuracy dự đoán đáp án cuối đúng/sai đạt 65.4\%, cao hơn 7.8 điểm phần trăm so với Logistic Regression theo độ dài (57.6\%) và 2.5 điểm phần trăm so với tree baseline LCoT2Tree (62.9\%) ở Bảng~1. Trong Best-of-$N$ ($N=5$) trên 600 câu MMLU-Pro của \texttt{DS-Qw7B}, selection bằng CoT2Graph đạt 54.5\%, so với ORM-Best 50.8\%, ORM + uncertainty 51.6\% và mean baseline 43.2\%; kết hợp entropy đạt 56.1\%.

Feng và cộng sự~\cite{rra3_effective_reasoning} so sánh mười reasoning model và xác định \emph{Failed-Step Fraction} (FSF) dự đoán correctness nhất quán hơn length hoặc review ratio. Selection theo FSF và can thiệp loại nhánh bị bỏ dở hỗ trợ control dựa trên cấu trúc. Tuy nhiên, graph phụ thuộc judge, và một nhánh bị bỏ dở không nhất thiết chứa bước sai logic.

\subsubsection{Giới hạn scaling theo độ phức tạp tổ hợp}
\textbf{ZebraLogic} \cite{rra3_zebralogic} đánh giá suy luận suy diễn qua câu đố lưới Constraint Satisfaction Problem (CSP), với search space từ $<10^3$ đến $>10^9$ ứng viên. Tác giả ghi nhận \emph{Curse of Complexity}: accuracy giảm mạnh khi search space vượt $10^6$. Parameter scaling không vượt được rào cản này: \texttt{Llama-3.1-405B} đạt 32.6\% tổng thể nhưng còn 1.5\% ở Large và 0.0\% ở X-Large. Reasoning model tốt hơn ở độ khó vừa vẫn suy giảm mạnh ở X-Large: \texttt{o1-full} 42.5\%, \texttt{DeepSeek-R1} 28.5\%, \texttt{o1-mini} 12.0\%. Repeated sampling (Best-of-$N$) và self-refinement chỉ đem lại cải thiện hạn chế khi độ phức tạp tổ hợp cao.

\subsubsection{Overthinking, underthinking và self-doubt}
Chen và cộng sự~\cite{rra3_chen_overthinking} đề xuất outcome/process efficiency metric và self-training để giảm overthinking. Trên MATH500, biến thể QwQ-32B-Preview giảm từ 2407.9 xuống 1330.7 token, trong khi accuracy đổi từ 93.0\% thành 92.8\%; tiết kiệm vì vậy có thể đi kèm mất hiệu năng. Ngược lại, Wang và cộng sự~\cite{rra3_underthinking} chẩn đoán việc đổi liên tục giữa các hướng chưa được phát triển đủ sâu. Trên AIME 2024, TIP tăng Pass@1 của QwQ-32B-Preview từ 38.3\% lên 44.1\%, ước lượng từ 32 mẫu mỗi bài. Tổng trace dài vẫn có thể đi cùng độ sâu exploration không đủ.

\textbf{Peng và cộng sự} \cite{rra3_overthinking} phân tích overthinking qua hành vi \emph{self-doubt}, tức kiểm tra lặp lại các kết luận trung gian đã đúng, như một nguồn tiêu token dư thừa quan trọng. Trên bài thiếu tiền đề, model không được hướng dẫn có thể sinh 4,700--6,400 token để cố chứng minh điều không thể. Prompt kiểm tra tiền đề giúp giảm 35\%--45\% token trong khi giữ hoặc tăng accuracy ở các model được thử.

% -------------------------------------------------------------------------
\subsection{Evaluation/Benchmark (A4)}
\label{subsec:a4_evaluation_benchmark}

Đánh giá Long-CoT cần benchmark kiểm tra state tracking qua horizon dài, metric tự động đo reasoning efficiency và giao thức process supervision.

\subsubsection{Benchmark về dependency dài và giới hạn scaling}
Để bổ sung cho benchmark ngắn đã bão hòa, \textbf{LongCoT} \cite{rra4_longcot} cung cấp 2,500 bài sinh từ template tham số hóa do chuyên gia thiết kế, thuộc năm miền: toán, hóa, khoa học máy tính, cờ vua và logic hình thức. Độ dài lời giải thay đổi lớn, thường 10k--100k+ token. Trên 2,000 bài vừa/khó, frontier model được thử có accuracy thấp: \texttt{GPT-5.2} đạt 9.83\%, \texttt{Gemini~3~Pro} 6.08\%, \texttt{Grok~4.1} 2.04\%. Trong thí nghiệm riêng tăng dependency DAG tuần tự từ 1 đến 35 node bài toán con, accuracy suy giảm nhanh hơn baseline lỗi độc lập ($p^N$), đặc biệt sau 15 node. Kết quả hỗ trợ giả thuyết của tác giả về sự suy giảm khả năng mang trạng thái qua context kéo dài.

\textbf{SPaRC} \cite{rra4_sparc} đánh giá spatial planning nhiều bước qua câu đố tìm đường lưới 2D biểu diễn bằng văn bản, chia năm mức khó. Người giải đạt 98.0\% tổng thể và 94.5\% ở Level~5; reasoning model cao nhất được thử, \texttt{o4-mini}, đạt 15.8\% tổng thể và 1.1\% ở Level~5, với hơn 50\% đường đi vi phạm luật di chuyển. SPaRC còn ghi nhận thiếu adaptive compute scaling: con người tăng thời gian suy nghĩ theo độ khó, trong khi \texttt{QwQ-32B} sinh lượng token gần nhau giữa các mức, 14,433 ở Level~1 và 13,114 ở Level~5.

\subsubsection{Khung tự động đánh giá redundancy dựa trên graph}
\textbf{CoTJudger} \cite{rra4_cotjudger} đánh giá efficiency không cần nhãn người từng bước bằng cách chuyển trace thành directed dependency graph và tìm \emph{Shortest Effective Path} (SEP). Phương pháp dùng forward depth-first search kết hợp LLM consistency judge. Redundancy Ratio được định nghĩa $R = (|V| - L_{\text{eff}})/|V|$, biểu thị tỷ lệ node/bước nằm ngoài SEP. Trên 21 reasoning model, nhiều node nằm ngoài SEP, tập trung trong các vòng self-verification sau khi đã có thể kết luận. SEP được xem là structural proxy phụ thuộc judge, không chứng minh mọi bước bị loại đều thừa về ngữ nghĩa. Tác giả cũng quan sát distilled model có logical density trên mỗi bước thấp hơn teacher.

\subsubsection{Benchmark process supervision và critique ở cấp bước}
\textbf{ProcessBench} \cite{rra4_processbench} cung cấp 3,400 trường hợp với lời giải từng bước trên bốn subset GSM8K, MATH, OlympiadBench và Omni-MATH. Chuyên gia gán nhãn bước lỗi đầu tiên. So sánh Process Reward Model (PRM) chuyên biệt với LLM critic cho thấy PRM học trên dữ liệu đơn giản tổng quát hóa kém ở bài đa dạng: \texttt{Skywork-PRM-7B} đạt F1 trung bình 42.1\% qua bốn subset. Reasoning model được prompt làm critic có kết quả cao hơn: \texttt{QwQ-32B-Preview} đạt F1 trung bình 71.5\%, vượt \texttt{GPT-4o} (61.9\%) nhưng thấp hơn \texttt{o1-mini} (87.9\%).

\textbf{CriticBench} \cite{rra4_criticbench} đánh giá ba năng lực Generation, Critique và Correction (GQC) trên năm miền suy luận. Trên 17 LLM, accuracy critique tương quan dương với khả năng correction sau đó; các model nhỏ có xu hướng dè dặt khi self-critique nếu không được hướng dẫn bằng external feedback.
